% =============================================================================
\subsubsection{Problem Setup}
% =============================================================================

We consider the fixed single-layer convolutional neural network
\begin{equation}
  f_\theta(x) = h(Ax),
\end{equation}
where $A\in\mathbb{R}^{N\times N}$ is a convolution matrix corresponding to a
fixed kernel $\theta$, and $h(\cdot)$ is an element-wise sigmoid activation:
\begin{equation}
  h(z)=\frac{1}{1+e^{-\alpha(z-c)}},\qquad
  h'(z)=\alpha\,h(z)\bigl(1-h(z)\bigr),
\end{equation}
where $\alpha>0$ is the stiffness parameter and $c>0$ is a threshold. The
input is constrained to $0\leq x\leq 1$, the target is binary
$y\in\{0,1\}^{N}$, and the reconstruction objective is
\begin{equation}
  J(x)=\|y-f_\theta(x)\|_2^2,\qquad
  x^*=\argmin_{0\leq x\leq 1}J(x).
  \label{eq:obj}
\end{equation}

We define the \textbf{gradient gate} $\Gamma(x)=h'(Ax)\in\mathbb{R}^N$, with
$\Gamma_i(x)=\alpha[f_\theta(x)]_i(1-[f_\theta(x)]_i)$. Coordinate $i$ is
\emph{active} if $|\Gamma_i(x)|>\varepsilon$. The \textbf{active gradient
fraction} is $F_\alpha(x)=|\{i:|\Gamma_i(x)|>\varepsilon\}|/N$.

For the two-layer extension $f_2(x)=h(A_2\cdot h(A_1 x))$, the gradient via
chain rule is
$\nabla J(x)=2A_1^\top[\Gamma^{(1)}\odot A_2^\top(\Gamma^{(2)}\odot(h_2-y))]$,
with effective gate
$\Gamma_{\mathrm{eff}}=\Gamma^{(1)}\odot A_2^\top\Gamma^{(2)}$---the source
of compounding collapse.

% =============================================================================
\subsubsection{Landscape Analysis}
% =============================================================================

The gradient is
\begin{equation}
  \nabla J(x)=2A^\top\bigl((h(Ax)-y)\odot h'(Ax)\bigr).
\end{equation}

\begin{lemma}[Gradient Bound]
\label{lem:grad_bound}
The magnitude of the gradient $\nabla J(x)$ is bounded by the stiffness
$\alpha$:
\begin{equation}
  \|\nabla J(x)\|_2\leq\frac{\alpha}{2}\|A^\top(h(Ax)-y)\|_2.
\end{equation}
\end{lemma}
\begin{proof}
Since $h'(z)=\alpha h(z)(1-h(z))\leq\alpha/4$, applying the triangle
inequality to the gradient expression yields the result.
\end{proof}

\begin{theorem}[Gradient Gate Bound]
\label{thm:gate}
For any $\varepsilon\in(0,\alpha/4)$ and $x\in[0,1]^N$:
\begin{equation}
  \mathbb{E}_x[F_\alpha(x)]\;\leq\;1-\frac{2\varepsilon}{\alpha}.
\end{equation}
As $\alpha\to\infty$ with $\varepsilon$ fixed, $F_\alpha(x)\to 0$ for
almost every $x$ with $(Ax)_i\neq c$ for all $i$.
\end{theorem}
\begin{proof}
The set $\{\Gamma_i>\varepsilon\}$ maps to a sigmoid-output interval
$(s_-,s_+)$ with
$s_\pm=\tfrac{1}{2}\pm\tfrac{1}{2}\sqrt{1-4\varepsilon/\alpha}$. Its width
satisfies $s_+-s_-=\sqrt{1-4\varepsilon/\alpha}\leq 1-2\varepsilon/\alpha$
(using $\sqrt{1-t}\leq 1-t/2$). When $(Ax)_i$ is approximately uniform,
$P(\Gamma_i>\varepsilon)\leq 1-2\varepsilon/\alpha$. As $\alpha\to\infty$,
the corresponding $z$-interval shrinks to $\{c\}$, giving $F_\alpha(x)\to 0$
for any $x$ not satisfying $(Ax)_i=c$ for all $i$.
\end{proof}

\begin{theorem}[Compounding Collapse --- Scoped Validity]
\label{thm:compound}
Under an approximate gate-independence assumption at convergence
(empirically verified: mean $|\mathrm{corr}(\Gamma^{(1)},A_2^\top\Gamma^{(2)})|=0.201$,
max $<0.50$ at all tested configurations), for the two-layer model with equal
stiffness $\alpha$ and kernel $A_2$:
\begin{equation}
  \frac{F^{(2)}_\alpha}{F^{(1)}_\alpha}
  \;\leq\;\frac{1}{1+c_2\alpha},\qquad
  c_2=\frac{\varepsilon}{\|A_2\|_\infty}.
\end{equation}
For $\alpha\geq16$, bootstrap $R^2\geq0.919$ (CI lower bound: 0.764). The
rate $c$ is not a universal constant; in the valid regime it follows
$c(\alpha)\approx0.211\cdot\alpha^{-0.674}$ ($R^2=0.973$), ranging from
$c\approx0.031$ at $\alpha=16$ to $c\approx0.012$ at $\alpha=60$. Below
$\alpha=16$, gradients are insufficiently sparse for the multiplicative
independence assumption to hold. \textbf{This asymptotic scope
($\alpha\geq16$, and an independence assumption verified only weakly even
in that regime---mean $|\mathrm{corr}|=0.201$, not strictly zero) is the
reason this theorem is not presented as a general law in the main paper.}
\end{theorem}
\begin{proof}[Proof sketch]
Coordinate $i$ is active in the two-layer model only if active in both
$\Gamma^{(1)}$ and $A_2^\top\Gamma^{(2)}$. Under approximate independence,
$P(\text{active in 2L})=P(\Gamma^{(1)}_i>\varepsilon)\cdot
P((A_2^\top\Gamma^{(2)})_i>\varepsilon)$. Applying Theorem~\ref{thm:gate}
to the second layer with effective threshold $\varepsilon/\|A_2\|_\infty$
gives the bound.
\end{proof}

\begin{lemma}[Adam Advantage from Momentum]
\label{lem:adam}
In the high-$\alpha$ regime, Adam's first moment
$m_t=\beta_1 m_{t-1}+(1-\beta_1)g_t$ preserves gradient signal from past
active sets via exponential decay at rate $\beta_1^k$. Oracle coordinate
rescaling ($\beta_1=0$) does not recover this advantage.

\noindent\emph{Verification (20 seeds):} Mean IoU --- Adam: $0.950$,
Oracle: $0.794$, PGD: $0.811$. Wilcoxon $p=0.0078$, Cohen's $d>9$ at all
$\alpha\geq5$. The Adam--Oracle gap is non-monotone, peaking at $0.332$ at
$\alpha=13$ with a local minimum of $0.246$ at $\alpha=25$.
\end{lemma}

\begin{figure}[htbp]
\centering
\includegraphics[width=0.45\linewidth]{landscape_alpha_2.png}
\hfill
\includegraphics[width=0.45\linewidth]{landscape_alpha_20.png}
\caption{Two-dimensional landscape slices near the solution. Low $\alpha$
(left) yields a smooth, well-conditioned surface. High $\alpha$ (right)
introduces large flat plateaus, sharp transitions, and loss of gradient
signal. The landscape changes qualitatively, not merely quantitatively, as
$\alpha$ increases.}
\label{fig:landscape}
\end{figure}

\paragraph{Identifiability and Null Space Geometry.}
\label{sec:null}

Reconstruction is unique if and only if the effective Jacobian
$J(x^*)=\mathrm{diag}(\Gamma(x^*))\cdot A$ has trivial null space,
$\sigma_{\min}(J(x^*))>0$.

\begin{figure}[htbp]
\centering
\includegraphics[width=0.95\linewidth]{nullspace_geometry.png}
\caption{Null space geometry. (Left) Null space dimension grows monotonically
with $\alpha$. (Centre) Minimum singular value $\sigma_{\min}(J(x^*))$
decreases monotonically, providing a constructive identifiability certificate.
(Right) Stable rank and condition number.}
\label{fig:nullspace}
\end{figure}

$\sigma_{\min}(J(x^*))$ decreases monotonically with $\alpha$, crossing zero
near $\alpha^*$. At $\alpha\geq20$, the null space dimension reaches $46/256$
(18\% of input space). The stable rank collapses monotonically from 46.8 at
$\alpha=1$ to 6.9 at $\alpha=60$ ($6.7\times$), computed via exact SVD of the
effective Jacobian, confirming the implicit dimensional reduction induced by
the gradient gate.

\paragraph{Uncertainty Quantification for $\alpha^*$.}
\label{sec:fisher}

To avoid parameter collinearity in the Fisher information matrix
($\mathrm{cond}(\mathcal{I}_4)=4.02\times10^6$ for the 4-parameter fit,
flagged numerically unreliable due to IoU$_{\max}$--$\delta$ collinearity,
$r=0.781$), we fix $\mathrm{IoU}_{\max}=1.0$ (justified by the data: Adam
achieves $\mathrm{IoU}=1.000$ at all $\alpha\leq10$) and fit a 3-parameter
model:
\begin{equation}
  \overline{\mathrm{IoU}}(\alpha)
  =\mathrm{IoU}_{\min}+
  \frac{1-\mathrm{IoU}_{\min}}
       {1+\exp((\alpha-\alpha^*)/\delta)}.
\end{equation}
This reduces $\mathrm{cond}(\mathcal{I}_3)$ to $3.0\times10^4$ ($135\times$
improvement). The primary uncertainty estimate is the distribution-free
bootstrap 95\% CI: $\alpha^*\in[27.92,\,29.78]$ (1000 resamples, 20 seeds),
corresponding to a CI width of 1.86. The 3-parameter CRLB gives
$\sigma_{\mathrm{CRLB}}(\alpha^*)=0.344$.

\begin{figure}[htbp]
\centering
\includegraphics[width=0.92\linewidth]{fisher_cramer_rao.png}
\caption{Uncertainty quantification for $\alpha^*$. (Left) 3-parameter
sigmoid fit with bootstrap CI $[27.92,29.78]$. The 4-parameter CRLB is
retired ($\mathrm{cond}(\mathcal{I})=4.02\times10^6$, unreliable); the
3-parameter replacement ($\mathrm{cond}=3.0\times10^4$) gives
$\sigma_{\mathrm{CRLB}}=0.344$. (Right) 3-parameter CRLB correlation matrix,
showing residual $\alpha^*$--$\delta$ correlation ($r=0.863$), structural for
this parameterisation but manageable.}
\label{fig:fisher}
\end{figure}

% =============================================================================
\subsubsection{Algorithmic Analysis}
% =============================================================================

\paragraph{Projected Gradient Descent (PGD).}
Relies directly on gradient magnitude; struggles severely when gradients
vanish. Applies uniform updates across all coordinates, failing to exploit the
few regions where gradients remain informative.

\paragraph{Momentum and Nesterov.}
Improve upon PGD by accumulating gradients over time, partially compensating
for sparse updates. Nesterov's look-ahead correction yields faster convergence
in smooth regimes, but the look-ahead step lands in an equally gradient-sparse
region at high $\alpha$.

\paragraph{Adam.}
Rescales gradients coordinate-wise based on historical gradient magnitudes.
Lemma~\ref{lem:adam} shows that Adam's advantage arises from momentum
accumulation, not coordinate-wise rescaling as commonly assumed.

\paragraph{L-BFGS-B.}
Approximates curvature information and can accelerate convergence in smooth
regimes. Performance degrades when gradients are sparse or unreliable.

\begin{table}[H]
\centering
\caption{Qualitative and quantitative comparison of optimisers. IoU and active
gradient fraction $F_\alpha$ at $\alpha=10$ and $\alpha=20$
(\texttt{sobel\_x}, checkerboard, 5 seeds).}
\label{tab:optimizer_summary}
\begin{tabular}{lcccccc}
\toprule
\multirow{2}{*}{Optimiser} & \multirow{2}{*}{Robustness} &
\multirow{2}{*}{Stability} &
\multicolumn{2}{c}{$\alpha=10$} &
\multicolumn{2}{c}{$\alpha=20$}\\
\cmidrule(lr){4-5}\cmidrule(lr){6-7}
 & & & IoU & $F_\alpha$ & IoU & $F_\alpha$\\
\midrule
\textbf{Adam} & High   & High     & \textbf{1.000} & \textbf{1.000} & \textbf{.947} & \textbf{.470}\\
L-BFGS-B      & Medium & Variable & .999 & .998 & .883 & .359\\
Momentum      & Medium & Medium   & 1.000 & 1.000 & .784 & .328\\
Nesterov      & Medium & Medium   & 1.000 & 1.000 & .754 & .278\\
PGD           & Low    & Low      & .858 & .896 & .787 & .270\\
\bottomrule
\end{tabular}
\end{table}

% =============================================================================
\subsubsection{Experiments}
% =============================================================================

All experiments use $64{\times}64$ inputs with five fixed convolutional
kernels: \textbf{identity-like} (trivial, separable), \textbf{avg\_blur}
($3{\times}3$ uniform, low-pass), \textbf{random\_norm} (random $3{\times}3$,
unit Frobenius), \textbf{sobel\_x} (horizontal edges, spectral norm
$\approx4.0$), and \textbf{laplacian} (second-order, spectral norm
$\approx8.0$). Threshold $c=0.5$; stiffness $\alpha\in\{1,2,5,10,20,40\}$
(core experiments) and $\alpha\in\{1,1.5,2,3,5,7,10,13,16,20,25,30,40,60\}$
(phase transition characterisation, 20 seeds). Adam uses $\mathrm{lr}=0.03$,
$\beta_1=0.9$, $\beta_2=0.999$; PGD and Momentum use $\mathrm{lr}=0.1$ with
$\beta=0.9$; Nesterov uses the same; L-BFGS-B uses its default line search.
All methods run for 200 iterations.

\paragraph{Metrics.}
Final loss $J(x)$; Intersection-over-Union (IoU, binarised at 0.5); active
gradient fraction $F_\alpha$ ($\varepsilon=0.01$); relative loss reduction
$(J(x_0)-J(x_T))/J(x_0)$; success rate $P[\mathrm{IoU}>0.7]$.

\paragraph{Phase Transition Characterisation.}
\label{sec:phase}

Mean reconstruction IoU follows a sigmoidal decay in $\alpha$ across a
transition band $\alpha\in[13,30]$. Three threshold values characterise the
band:
\begin{center}
\begin{tabular}{lcc}
\toprule
Threshold & $\alpha$ & IoU at crossing\\
\midrule
$\alpha_{10}$: IoU first drops below 1.000 & 13 & 0.996\\
$\alpha_{50}$: IoU drops below 0.900 & 30 & 0.875\\
$\alpha_{90}$: IoU drops below 0.800 & $\geq60$ & 0.771\\
\bottomrule
\end{tabular}
\end{center}

The 3-parameter sigmoid fit ($\mathrm{IoU}_{\max}$ fixed at 1.0) gives:
\[
  \mathrm{IoU}_{\min}=0.782,\quad
  \alpha^*=28.97,\quad
  \delta=6.29.
\]
The bootstrap median is $\alpha^*=28.83$ (95\% CI $[27.92,\,29.78]$).
Earlier drafts reported $\alpha^*\approx11.7$ as an artefact of a bounded
least-squares fit applied to a 6-point sparse grid whose densest coverage
lay entirely within the subcritical plateau ($\mathrm{IoU}\equiv1.0$); the
bound $\alpha^*\leq11.7$ was itself drawn from a preliminary estimate,
making the result circular. With a 14-point grid extending to $\alpha=60$
(20 independent seeds, 1000-resample bootstrap), the full sigmoid S-curve
is observed and the above estimate is obtained.

\begin{figure}[htbp]
\centering
\includegraphics[width=0.72\linewidth]{phase_transition_fit.png}
\caption{Phase transition fit. IoU follows a sigmoidal curve in $\alpha$ with
bootstrap midpoint $\alpha^*=28.83$ (95\% CI $[27.92,29.78]$, dashed line
and shaded band). Transition onset ($\alpha_{10}=13$) and degradation
midpoint ($\alpha_{50}=30$) are marked. Networks with $\alpha\leq10$ achieve
perfect reconstruction ($\mathrm{IoU}=1.000$); quality degrades continuously
through the band $\alpha\in[13,30]$.}
\label{fig:phase_fit}
\end{figure}

\paragraph{Effect of Sigmoid Stiffness.}

\begin{figure}[htbp]
\centering
\includegraphics[width=0.47\linewidth]{alpha_sweep_loss_mean_std.png}
\hfill
\includegraphics[width=0.47\linewidth]{alpha_sweep_iou_mean_std.png}
\caption{Effect of sigmoid stiffness on final loss (left) and reconstruction
IoU (right) across all optimisers. Adam achieves $\mathrm{IoU}=1.000$ at
$\alpha\leq10$ and degrades gracefully to $0.830$ at $\alpha=40$.}
\label{fig:alpha_main}
\end{figure}

Performance degrades consistently as $\alpha$ increases through the transition
band (Fig.~\ref{fig:alpha_main}). Adam sustains perfect reconstruction
($\mathrm{IoU}=1.000$) up to $\alpha=10$, then degrades smoothly to 0.947 at
$\alpha=20$ and 0.830 at $\alpha=40$. The non-adaptive optimisers begin
degrading earlier, with PGD already at IoU $=0.858$ at $\alpha=10$.

\begin{table}[h]
\centering
\caption{Performance across sigmoid stiffness values (Adam, \texttt{sobel\_x},
20 seeds).}
\label{tab:alpha_adam}
\begin{tabular}{ccc}
\toprule
$\alpha$ & Loss (mean$\pm$std) & IoU (mean$\pm$std)\\
\midrule
1.0  & $504.87\pm9.91$   & $1.000\pm0.000$\\
5.0  & $344.19\pm11.62$  & $1.000\pm0.000$\\
10.0 & $488.05\pm9.10$   & $1.000\pm0.000$\\
20.0 & $646.31\pm11.35$  & $0.947\pm0.007$\\
40.0 & $799.49\pm30.10$  & $0.830\pm0.012$\\
60.0 & $892.13\pm34.52$  & $0.771\pm0.018$\\
\bottomrule
\end{tabular}
\end{table}

\paragraph{Impact of Cutoff Threshold $c$.}

\begin{figure}[htbp]
\centering
\includegraphics[width=0.72\linewidth]{threshold_sensitivity_analysis.png}
\caption{Sensitivity analysis of the cutoff threshold $c$. Performance is
maximised at $c=0.5$ and degrades symmetrically toward the boundaries. Adam
shows greater resilience to off-centre thresholds than PGD or L-BFGS-B.}
\label{fig:threshold_sweep}
\end{figure}

\paragraph{Spatial Gradient Gate and Dimensional Collapse.}

The gradient gate $\Gamma(x)=h'(Ax)$ collapses spatially as $\alpha$
increases: a dense, near-uniform field at low $\alpha$ narrows to isolated
active coordinates at edge boundaries at high $\alpha$, directly
instantiating Theorem~\ref{thm:gate}. The quantitative version of this
collapse, from \texttt{grad\_sparsity\_all\_optimizers.csv}: at $\alpha=1$, the gate is 100\% active. The collapse becomes pronounced in
the transition band: by $\alpha=20$, Adam retains 47.0\% active coordinates
versus PGD's 27.0\%; by $\alpha=40$, Adam retains 19.0\% versus PGD's 5.7\%.
Stable rank collapses monotonically from 46.8 at $\alpha=1$ to 6.9 at
$\alpha=60$ ($6.7\times$), computed via exact SVD of the effective Jacobian.

\paragraph{Gradient Sparsity and Its Impact.}

\begin{figure}[htbp]
\centering
\includegraphics[width=0.65\linewidth]{fig_gradient_sparsity_per_optimizer.png}
\caption{Fraction of active gradient components ($|\Gamma_i|>0.01$) at
convergence vs.\ $\alpha$. Adam maintains 47.0\% active at $\alpha=20$
compared to PGD's 27.0\%, and 19.0\% vs.\ 5.7\% at $\alpha=40$---a widening
gap deep in the transition band that explains Adam's robustness.}
\label{fig:grad_activity}
\end{figure}

\FloatBarrier

\paragraph{Oracle Ablation: Momentum Drives Adam's Advantage.}
\label{sec:oracle}

\begin{figure}[htbp]
\centering
\includegraphics[width=0.92\linewidth]{oracle_ablation.png}
\caption{Oracle ablation: Adam vs.\ Oracle Rescaling vs.\ PGD. (Left) IoU
curves showing Oracle $<$ PGD at $\alpha=10$--$25$ (shaded zone) and Adam's
sustained superiority ($p=0.0078$, all $\alpha\geq5$). (Right) Adam--Oracle
gap peaks at 0.332 ($\alpha=13$) and anticorrelates with $|d\Gamma/d\alpha|$
($r=-0.799$, $p=0.0006$), revealing the transition-regime miscalibration
mechanism.}
\label{fig:oracle}
\end{figure}

We construct Oracle Rescaling: PGD with per-coordinate step sizes $v_t^{-1/2}$
from Adam's second moment, but $\beta_1=0$ (no momentum). Mean IoU: Adam
$0.950$, Oracle $0.794$, PGD $0.811$ ($p=0.0078$, Wilcoxon, Cohen's
$d>9$, at all $\alpha\geq5$).

\paragraph{Reversal of Oracle vs.\ PGD.}
Oracle performs \emph{below} plain PGD at $\alpha=10$--$25$ (gap $-0.093$ to
$-0.183$ IoU points, $p=0.0078$--$0.023$). This reversal is not explained by
differential gate collapse: at these $\alpha$ values, Oracle retains 12--63\%
more active gradient coordinates than PGD (Oracle/PGD active-fraction ratio:
1.12 at $\alpha=10$, rising to 1.64 at $\alpha=25$), yet achieves
substantially lower IoU. Instead, the reversal tracks the rate of change of
gate sparsity: the Oracle--PGD gap anticorrelates strongly with
$|d\Gamma/d\alpha|_{\rm PGD}$ (Pearson $r=-0.799$, $p=0.0006$), with the
deepest reversal at $\alpha=13$--$16$ where gate density falls most steeply.
We attribute this to miscalibration of exact second-moment rescaling in a
rapidly-changing gradient field: oracle step sizes assign disproportionately
large updates to barely-active coordinates (low $v_t\to$ high $1/\sqrt{v_t}$)
near the sigmoid inflection point, pulling the reconstruction trajectory away
from the true attractor. PGD's fixed step implicitly weights updates by
gradient magnitude, which is already calibrated by the gate structure itself.
The reversal resolves above $\alpha=30$, where gate density has settled to a
new stable regime ($|d\Gamma/d\alpha|\approx0$).

\begin{quote}
\emph{Practical implication:} The advantage is in the gradient history, not
the step-size schedule. Tuning per-coordinate learning rates does not replicate
Adam's advantage in high-$\alpha$ regimes.
\end{quote}

\paragraph{Compounding Collapse: Theorem~\ref{thm:compound} Verification.}

\begin{figure}[htbp]
\centering
\includegraphics[width=0.92\linewidth]{depth_scaling_law.png}
\caption{Depth scaling law. (Left) Bootstrap $R^2$ vs.\ $\alpha$: model is
valid ($R^2\geq0.919$) only for $\alpha\geq16$; $R^2$ is negative for
$\alpha\leq13$ (model inapplicable, marked $\times$). (Centre) Collapse rate
$c(\alpha)\approx0.211\cdot\alpha^{-0.674}$ decreases monotonically in the
valid regime. (Right) Log-linear collapse confirmed for $\alpha\geq16$.}
\label{fig:depth}
\end{figure}

\begin{table}[H]
\centering
\caption{Compounding collapse: empirical ratio (2L/1L) vs.\ bound
$1/(1+c(\alpha)\cdot\alpha)$. Valid regime: $\alpha\geq16$ ($R^2\geq0.919$).}
\label{tab:compounding}
\begin{tabular}{ccccccc}
\toprule
$\alpha$ & 1L $F$ & 2L $F$ & Ratio & $c_{\rm fit}$ & $R^2$ (boot) & Applicable\\
\midrule
1--13 & --- & --- & --- & --- & $<0$ & \textbf{No}\\
16 & --- & --- & --- & 0.031 & 0.919 & Yes\\
20 & 0.147 & 0.078 & 0.531 & 0.029 & 0.938 & Yes\\
25 & --- & --- & --- & 0.026 & 0.959 & Yes\\
40 & 0.059 & 0.018 & 0.302 & 0.017 & 0.972 & Yes\\
60 & --- & --- & --- & 0.012 & 0.977 & Yes\\
\bottomrule
\end{tabular}
\end{table}

The gate independence assumption is empirically verified: mean
$|\mathrm{corr}(\Gamma^{(1)},A_2^\top\Gamma^{(2)})|=0.201$ with no seed
exceeding 0.50. This introduces a small positive bias in the predicted
compounding rate.

\paragraph{Activation Comparison.}

\begin{figure}[htbp]
\centering
\includegraphics[width=0.95\linewidth]{activation_taxonomy_v2.png}
\caption{Activation comparison across six functions. (Left) IoU vs.\ $\alpha$.
(Centre) Active gradient fraction---three regimes are clearly separated.
(Right) Silhouette score vs.\ $k$: $k=3$ is optimal (silhouette$=0.722$),
confirming a 3-group taxonomy. Between-group variance is $573\times$
within-group variance at $\alpha=40$.}
\label{fig:activation}
\end{figure}

Six activations reveal a three-group taxonomy (Fig.~\ref{fig:activation},
$k=3$, silhouette$=0.722$):
\begin{itemize}[topsep=2pt,itemsep=1pt]
\item \textit{Saturating} (sigmoid, tanh): active fraction collapses
  $\approx60\%$ from $\alpha=20$ to $\alpha=40$ (from 0.459 to 0.182);
  IoU degrades proportionally. Note: tanh degrades more severely in IoU
  ($0.598$ at $\alpha=40$) compared to sigmoid ($0.830$), attributable to
  tanh's bipolar output range $[-1,1]$ introducing additional optimisation
  difficulty with binary targets.
\item \textit{Smooth self-gated} (GELU, Mish, SiLU/Swish): active fraction
  remains $>95\%$ at all tested $\alpha$---no collapse regardless of
  stiffness, within \emph{this} synthetic fixed-kernel reconstruction
  setting. This is the same qualitative direction the main paper finds
  during ordinary training of real networks (Section~\ref{sec:dynamics}),
  though the two settings are connected empirically
  (Section~\ref{sec:smoothness}), not assumed identical.
\item \textit{Piecewise linear} (ReLU): active fraction is approximately
  $\alpha$-insensitive; IoU degradation is moderate and stable.
\end{itemize}

The prior label ``plateau'' for Swish was incorrect; Swish clusters with GELU
and Mish in the smooth group in all $k\geq2$ clusterings (silhouette-verified).

\paragraph{Kernel Sensitivity.}

\begin{figure}[htbp]
\centering
\includegraphics[width=0.47\linewidth]{kernel_sweep_loss_mean_std.png}
\hfill
\includegraphics[width=0.47\linewidth]{kernel_sweep_iou_mean_std.png}
\caption{Final loss (left) and IoU (right) across convolution kernels.
Identity-like and avg\_blur are always trivial ($\mathrm{IoU}=1.0$). Laplacian
produces the highest loss, consistent with its spectral norm
$\approx8.0$ inducing maximum landscape anisotropy.}
\label{fig:kernel}
\end{figure}

\begin{table}[H]
\centering
\caption{Kernel difficulty ranking: IoU (Adam, 5 seeds) and spectral norm.}
\label{tab:kernel_difficulty}
\begin{tabular}{lcccccc}
\toprule
Kernel & $\alpha=1$ & $\alpha=5$ & $\alpha=10$ & $\alpha=40$ &
$\|A\|_2$ & Difficulty\\
\midrule
identity-like & 1.000 & 1.000 & 1.000 & 1.000 & $\approx1.0$ & Trivial\\
avg\_blur     & 1.000 & 1.000 & 1.000 & 1.000 & $\approx1.0$ & Trivial\\
random\_norm  & 1.000 & 1.000 & 1.000 & 0.777 & $\approx1.0$ & Moderate\\
sobel\_x      & 1.000 & 1.000 & 1.000 & 0.830 & $\approx4.0$ & Hard\\
laplacian     & 1.000 & 0.964 & 0.812 & 0.589 & $\approx8.0$ & Very Hard\\
\bottomrule
\end{tabular}
\end{table}

\paragraph{Phase Diagram.}

\begin{figure}[htbp]
\centering
\includegraphics[width=0.76\linewidth]{phase_diagram_alpha_x_kernel.png}
\caption{Phase diagram: IoU heatmap over (kernel,$\alpha$) space. Spectral
norm and $\alpha$ jointly control difficulty.}
\label{fig:phase}
\end{figure}

\paragraph{Target Sensitivity.}

\begin{figure}[htbp]
\centering
\includegraphics[width=0.57\linewidth]{target_sweep_iou.png}
\caption{Reconstruction IoU across target patterns. Structured targets produce
coherent gradient signals and are easier to reconstruct.}
\label{fig:target}
\end{figure}

\paragraph{Generalisation and Universality.}
\label{sec:generalisation}

\begin{figure}[htbp]
\centering
\includegraphics[width=0.92\linewidth]{generalization_figure.png}
\caption{Generalisation results. Scale invariance: curves for
$32{\times}32$, $64{\times}64$, $128{\times}128$ nearly overlap. The
phase transition is a property of $(\alpha,\text{kernel})$, not input
resolution.}
\label{fig:gen}
\end{figure}

\textit{Scale invariance.}
IoU at $\alpha=10$: $32{\times}32$ (0.716), $64{\times}64$ (0.767),
$128{\times}128$ (0.786). Differences are within one standard deviation.

\textit{Finite-size scaling and universality.}
Two distinct transition thresholds are reported and should not be confused. The
finite-size scaling (FSS) estimate $\alpha^*_{\rm FSS}=17.2\pm0.6$ (64$\times$64)
is the stiffness at which the order parameter
$(\mathrm{IoU}-\mathrm{IoU}_{\min})$ first departs from its plateau maximum
(transition \emph{onset}, 81\% of maximum order parameter retained at this
point). The bootstrap estimate $\alpha^*=28.83$ (95\% CI $[27.92,29.78]$) is
the sigmoid \emph{midpoint}: the stiffness where
$\mathrm{IoU}=(\mathrm{IoU}_{\max}+\mathrm{IoU}_{\min})/2\approx0.893$. The
two quantities differ because the transition has finite width ($\delta\approx6.3$);
for a sharp transition ($\delta\to0$) they would coincide. The FSS estimate
characterises when reconstruction quality \emph{begins} to degrade; the sigmoid
midpoint characterises the \emph{centre} of the degradation regime.

\begin{figure}[htbp]
\centering
\includegraphics[width=0.92\linewidth]{finite_size_scaling.png}
\caption{Finite-size scaling. (Left) Critical exponents $\beta=0.030$--$0.043$
across resolutions, with correct error bars. The prior claim of $\beta=0.05$
was a boundary-clamping artefact (parameter pinned at lower bound with
$\pm57$--$839$ error). (Right) The two $\alpha^*$ definitions illustrated:
FSS onset ($\alpha^*_{\rm FSS}\approx17$) vs.\ bootstrap midpoint
($\alpha^*=28.83$), separated by approximately one transition width
$\delta\approx6.3$.}
\label{fig:fss}
\end{figure}

\textit{Noise robustness.}
Reconstruction quality varies $<6\%$ IoU across SNR 15--60\,dB for all
$\alpha$ (maximum spread 6.2\% at $\alpha=25$, within the transition band).
The dominant factor is $\alpha$, not noise level.

\paragraph{Convergence Behaviour.}

\begin{figure}[htbp]
\centering
\includegraphics[width=0.65\linewidth]{convergence_bands_loss.png}
\caption{Convergence curves with mean and variance bands. Adam demonstrates
both faster convergence and lower variance. Variability across runs increases
significantly with $\alpha$, confirming multiple local minima separated by
flat regions.}
\label{fig:convergence}
\end{figure}

\paragraph{Success Rate Analysis.}

\begin{figure}[htbp]
\centering
\includegraphics[width=0.65\linewidth]{optimizer_success_heatmap.png}
\caption{Success rate $P[\mathrm{IoU}>0.7]$ across optimisers and $\alpha$
values.}
\label{fig:success}
\end{figure}

\paragraph{Predictive Difficulty Model.}

\begin{figure}[htbp]
\centering
\includegraphics[width=0.95\linewidth]{predictive_difficulty_model.png}
\caption{Predictive difficulty model, fit from raw data and evaluated with
5-fold cross-validation. (Left) ROC curve, AUC$=0.958$. (Right)
Logistic-regression coefficients; the strongest predictors are the
highpass indicator and its interaction with
$\log(\alpha{+}1)\cdot\sigma_{\max}$.}
\label{fig:pred}
\end{figure}

A logistic regression over $(\log(\alpha+1)$, $\sigma_{\max}$,
$\mathbf{1}_{\mathrm{hp}}$, $\alpha\cdot\mathbf{1}_{\mathrm{hp}}$,
$\log(\alpha+1)\cdot\sigma_{\max})$, refit directly from
\texttt{phase\_diagram\_alpha\_x\_kernel.csv} ($n=560$) and evaluated with
5-fold cross-validation, predicts $P[\mathrm{IoU}>0.7]$ before running any
optimisation: \textbf{AUC$=0.958$}, \textbf{accuracy$=0.950$}.

\paragraph{Non-Uniqueness Analysis.}

\begin{figure}[htbp]
\centering
\begin{minipage}{0.48\linewidth}
  \includegraphics[width=\linewidth]{nonuniqueness_mosaic.png}
\end{minipage}
\hfill
\begin{minipage}{0.48\linewidth}
  \includegraphics[width=\linewidth]{nonuniqueness_input_vs_output_scatter.png}
\end{minipage}
\caption{Non-uniqueness at $\alpha=20$ from 20 random restarts. Large input
diversity maps to small output variation, directly quantifying the
near-degeneracy established in Section~\ref{sec:null}.}
\label{fig:nonunique}
\end{figure}

\paragraph{Privacy Implications: Gradient Leakage Attack Surface.}

\begin{figure}[htbp]
\centering
\includegraphics[width=0.92\linewidth]{gradient_leakage.png}
\caption{Gradient leakage attack surface (synthetic fixed-kernel setting).
Above $\alpha\approx14$, attack success rate falls below 20\%. The smooth
transition in the dense-$\alpha$ sweep confirms continuity. Compare to the
main paper's real-architecture, real-attack probe
(Section~\ref{sec:privacy}), which finds an attack-dependent rather than
uniformly protective effect.}
\label{fig:leakage}
\end{figure}

This synthetic setting is a minimal model of gradient inversion
attacks~\cite{zhu2019deep}. The transition band directly defines a
\emph{privacy boundary} within this model: operating above $\alpha\approx14$
(within the transition onset $\alpha_{10}=13$) provides strong natural
privacy protection at the cost of reduced model utility---a finding the
main paper's real-attack probe (Section~\ref{sec:privacy}) refines, finding
the protective effect itself is attack-dependent rather than uniform.

\paragraph{Reconstruction Under Learned Weights.}
\label{sec:learned}

\begin{figure}[htbp]
\centering
\includegraphics[width=0.95\linewidth]{learned_weights.png}
\caption{Fixed vs.\ jointly-optimised kernel reconstruction. (Left) IoU
comparison ($p=0.0078$ at all 14 $\alpha$ values, $n=8$). (Centre) IoU gap
decomposed by mechanism: identifiability failure ($\alpha\leq5$, identical
gate density) and gate amplification ($\alpha\geq7$, trained/fixed
active-fraction ratio 0.28--0.72). (Right) Trained kernel induces
effective stiffness 2--3$\times$ the nominal $\alpha$.}
\label{fig:learned}
\end{figure}

Trained networks are consistently harder to invert at every $\alpha$ value
($p=0.0078$, paired Wilcoxon, $n=8$, 14 $\alpha$ values). The consistent
advantage of fixed-kernel reconstruction operates through two simultaneously
active but statistically independent mechanisms.

At subcritical stiffness ($\alpha\leq5$), gate density is identical between
conditions (active fraction $=1.000$ for both), yet trained-kernel IoU is
$0.25$--$0.71$ versus fixed-kernel IoU of $1.000$. The gate framework does
not explain this gap: it reflects identifiability failure of joint
optimisation. When both the forward operator and input are free variables, the
optimiser can minimise reconstruction loss by adapting the kernel rather than
recovering the true input, making the problem ill-posed regardless of gate
structure.

At supercritical and transition-regime stiffness ($\alpha\geq7$), a second
mechanism activates: the jointly-optimised kernel induces gate collapse
2--3.5$\times$ deeper than the same nominal $\alpha$ produces under a fixed
kernel (trained/fixed active-fraction ratio: 0.54 at $\alpha=7$, 0.41 at
$\alpha=13$--$20$, 0.28 at $\alpha=60$). Joint optimisation effectively shifts
the system into a higher-stiffness regime, consistent with the prediction that
kernels with higher effective spectral norm produce steeper gate transitions
(Section~\ref{sec:null}, Table~\ref{tab:kernel_difficulty}).

The two pathways are statistically independent: the gate density deficit does
not predict the IoU gap magnitude across $\alpha$ values (Pearson $r=0.082$,
$p=0.82$). This finding advises against joint kernel--input optimisation in
gated reconstruction settings.

\paragraph{DeepNet Extension: ResNet-18 and VGG-11 (Synthetic Alpha-Sigmoid Variant).}
\label{sec:deepnet}

\begin{figure}[htbp]
\centering
\includegraphics[width=0.95\linewidth]{deepnet_gate_compounding.png}
\caption{DeepNet compounding collapse, using ResNet-18/VGG-11-\emph{shaped}
networks with every activation replaced by the synthetic alpha-sigmoid gate
(not the unmodified, real architectures studied in the main paper). (Left)
Bootstrap $R^2$ vs.\ $\alpha$: law is valid for $\alpha\geq20$ in both
architectures ($R^2=0.615$--$0.837$). (Centre) Collapse rates: VGG-11
($c\approx0.023$) matches the simple convnet; ResNet-18 ($c\approx0.009$)
is lower due to skip connections. (Right) Layer-wise collapse front
propagating from late to early layers.}
\label{fig:deepnet}
\end{figure}

\textbf{This subsection's ResNet-18/VGG-11 are alpha-sigmoid-shaped
synthetic networks, not the unmodified, canonical architectures studied in
the main paper's Sections~\ref{sec:dynamics}--\ref{sec:mechanism}.} Every
ReLU in these networks is replaced by the tunable alpha-sigmoid gate to
make a stiffness sweep meaningful, since canonical ReLU has no stiffness
parameter. Conclusions here are about gate compounding under heterogeneous
depth/width in this synthetic variant; the main paper's real-architecture
findings are established independently, with real, unmodified ReLU and
GELU/SiLU/Mish, and do not rely on this subsection.

The compounding collapse law extends to these alpha-sigmoid-shaped
ResNet-18 (17 layers) and VGG-11 (9 layers) variants. At $\alpha\leq10$,
$R^2$ is negative for both architectures---the same threshold as the
simple convnet. At $\alpha\geq20$: ResNet-18 $R^2 = 0.615$--$0.627$;
VGG-11 $R^2 = 0.774$--$0.837$.

Between $\alpha=10$ and $\alpha=20$, ResNet-18 collapses from 93\% active to
12\% ($7.7\times$ reduction); VGG-11 from 90\% to 12\% ($7.5\times$). A
collapse front propagates from late layers to early layers: at $\alpha=20$,
layers 11--16 in ResNet-18 are fully collapsed (active fraction $=0$) while
layers 0--6 retain 6--36\% activity. Skip connections locally sustain gradient
flow in residual blocks (layers 5--6: 31--32\% active at $\alpha=20$ vs.\
adjacent non-residual layers at 6--8\%).

The $\alpha\geq16$ validity threshold is architecture-independent within
this synthetic variant, suggesting it is a property of activation stiffness
rather than of depth, width, or skip connection density. Architecture-specific
collapse rates ($c_{\rm VGG}\approx0.023\approx c_{\rm convnet}$;
$c_{\rm ResNet}\approx0.009$) reflect structural differences in gradient
flow; a unified $c(\alpha,\mathrm{architecture})$ model is left for future
work.

\paragraph{Gate Independence Verification.}
\label{sec:independence}

\begin{figure}[htbp]
\centering
\includegraphics[width=0.85\linewidth]{theorem_verification.png}
\caption{Gate independence test for Theorem~\ref{thm:compound}. (Left) Mean
Pearson correlation between $\Gamma^{(1)}$ and $A_2^\top\Gamma^{(2)}$ at
convergence across $\alpha$ values. (Right) Distribution of correlations across
all seeds and $\alpha$ values. No seed exceeds $|\mathrm{corr}|=0.50$.}
\label{fig:gate_independence}
\end{figure}

The independence assumption underlying Theorem~\ref{thm:compound} is
empirically tested, not simply assumed. Mean $|\mathrm{corr}(\Gamma^{(1)},
A_2^\top\Gamma^{(2)})|=0.201$ across all tested configurations; no seed at
any $\alpha$ exceeds $|\mathrm{corr}|=0.50$. The gates are weakly correlated
rather than independent, introducing a small positive bias in the predicted
compounding rate, and---as stated at Theorem~\ref{thm:compound}---this is
why the theorem's validity is scoped to $\alpha\geq16$ rather than presented
as unconditional.

\paragraph{Discussion (Synthetic Setting).}

Three cross-cutting observations emerge from the synthetic fixed-kernel
setting. First, increasing $\alpha$ induces implicit dimensional collapse:
stable rank collapses $6.7\times$ from $\alpha=1$ (46.8) to $\alpha=60$
(6.9). Second, high-$\alpha$ regimes produce characteristic failure modes:
the optimiser stalls in flat plateaus, and sensitivity to initialisation
grows sharply (rising variance in Fig.~\ref{fig:convergence}). Third,
edge-driven dynamics: the gradient gate is active only when $Ax\approx c$,
corresponding to edge-like transition regions; early iterations reconstruct
coarse boundaries while later iterations fill interior regions.

Connection to compressed sensing: when $F_\alpha(x^*)\cdot N<\mathrm{rank}(A)$,
the effective system is underdetermined and recovery fails almost surely. This
suggests applying RIP-based guarantees from CS theory by treating the active
support as the effective measurement set~\cite{donoho2006compressed}.
